\documentclass[11pt,a4paper]{article}
\usepackage[utf8]{inputenc}
\usepackage[T1]{fontenc}
\usepackage{lmodern}
\usepackage[a4paper,margin=2.3cm]{geometry}
\usepackage{booktabs}
\usepackage{amsmath,amssymb}
\usepackage[colorlinks=true,linkcolor=blue,urlcolor=blue]{hyperref}
\setlength{\parskip}{0.55em}
\setlength{\parindent}{0pt}
\newcommand{\GS}{G\"ombSoma}
\title{Plan: \GS{} optimization passes 5--8\\
       \large{Targeted speedup of the remaining 28 ms per-image cost}}
\date{2026-05-15}
\author{}
\begin{document}
\maketitle

\section{Background}

Passes 1--4 took the \GS{} \texttt{Config E} per-image forward from
8\,283 ms to 28.0 ms (296$\times$). The remaining cost is spread
across many small components rather than a single dominant hot spot:

\begin{center}
\begin{tabular}{lrr}
\toprule
Component                       & ms    & \% of total \\
\midrule
StimulusGraphBuilder (×2 in Config E) & 9--10 & 33 \\
``Unaccounted''                  & 10--11 & 37 \\
Quadtree                         & 2.2   & 8 \\
SDRF                             & 3.0   & 11 \\
Bochner × 3 branches             & 1.4   & 5 \\
Patch encoder                    & 0.2   & 1 \\
\bottomrule
\end{tabular}
\end{center}

\section{Scope}

Four targeted optimisation passes, each addressing one slice of the
remaining cost. None touches CORE.YAML-protected files. Each is
independently revertable. The 150-test Ricci-Stim suite must pass
after every pass.

\subsection{P5 --- Profile the ``unaccounted'' 10 ms}

\textbf{Hypothesis:} sparse-tensor coalesce calls on M\_v constructions
plus Python dispatch overhead in the per-image batch loop.

\textbf{Fix:} use \texttt{torch.profiler} (kineto) to capture per-op
timing on a single forward pass at the Config E shape. Identify which
specific ops eat the 10 ms.

\textbf{Expected:} 10\,ms $\to$ 3\,ms (save 7\,ms) once the hot ops
are identified and replaced with cheaper equivalents.

\textbf{Risk:} the unaccounted may be many small ops with no single
fix. If so, fall back to a sub-ms profile and accept the 10 ms as
fundamental Python orchestration overhead, blocked on P8.

\subsection{P6 --- Cache topology between SDRF passes 1 and 2}

\textbf{Hypothesis:} \texttt{StimulusGraphBuilder} is called twice
in Config E. Both calls use the same quadtree, so:

\begin{itemize}
\item \texttt{\_same\_scale\_edges} produces byte-identical output
\item \texttt{\_cross\_scale\_edges} produces byte-identical output
\item adjacency dict is identical
\end{itemize}

Only the edge \emph{signs} and the SDRF-added shortcut edges
differ between calls.

\textbf{Fix:} add a \texttt{cached\_topology} parameter to
\texttt{StimulusGraphBuilder.forward}; pass 1 caches; pass 2 reuses
the cached topology and only computes signs + the augmented walks /
polygons / triangles incident to the new edges.

\textbf{Expected:} 9.5 ms $\to$ ~3 ms on the second pass. Total
Config E save: ~6 ms.

\textbf{Risk:} caching wrong objects (e.g., tensors with stale
gradients) would corrupt the forward. The cache key must be
\texttt{(id(tree), tree.positions.data\_ptr(), tree.sizes.data\_ptr())}
or equivalent.

\subsection{P7 --- Vectorise polygon enumeration}

\textbf{Hypothesis:} the polygon plaquette enum at $\sim$200 polygons
is still a Python loop over \texttt{(r, c, s)} tuples with a dict
lookup per side. The vectorised walk / triangle enum did not help
because (n, n, n) tensor allocation overhead matches the Python
cost at $n = 256$. Polygons are simpler: just check four neighbors
per anchor. The vectorisation pays off here because the per-anchor
work is structured (top-left $\to$ right $\to$ bottom $\to$ diag).

\textbf{Fix:} sort anchors by (scale, r, c) so adjacency is
stride-1; build a $(n_\text{anchors}, 4)$ neighbour-index tensor;
check the plaquette condition via tensor masks.

\textbf{Expected:} 3--4 ms saved inside StimulusGraphBuilder.

\textbf{Risk:} edge-index lookup for plaquette edges must match
the caller's edge ordering; mismatches caused a bug in P3
($\sigma$-product test caught it). Same lookup-via-shared-dict
pattern applies.

\subsection{P8 --- GPU-level batching across images}

\textbf{Hypothesis:} the per-image Python forward (\texttt{for b in
range(batch\_size): outs.append(\_forward\_single(images[b]))}) is
fundamentally Python-bound. To break out, we need a
SegmentedSparse representation where all images' anchors live in
one big sparse graph and per-image offsets index them.

\textbf{Fix:}
\begin{enumerate}
\item Concatenate all images' anchor positions / sizes into single
  tensors, with a per-image offset array.
\item Concatenate quadtree-derived edges with per-image offset
  shifts (so anchor indices in image B don't collide with image A).
\item Run one big \texttt{StimulusGraphBuilder} on the concatenated
  graph.
\item Run one big Bochner-coupled \texttt{HypergraphConv} on the
  concatenated graph.
\item Split outputs back per-image at the head.
\end{enumerate}

\textbf{Expected:} 5--10$\times$ on per-image-amortised cost. At
batch\_size = 8 this means 28 ms / image $\to$ 3--5 ms / image
amortised.

\textbf{Risk:} multi-session task. Touches every module's forward
signature. Plan as its own phase if the smaller fixes don't get
us into the desired range.

\section{Test strategy}

After each pass:

\begin{itemize}
\item Full Ricci-Stim 150-test suite passes.
\item Profile reruns: P5-target ms $\to$ measured ms shown in
  before/after table.
\item Smoke run (30 imgs $\times$ 1 epoch $\times$ Config E) wall
  time recorded; expected monotone decrease.
\end{itemize}

\section{Performance budget}

Per-pass GPU time $\leq$ 0 (these are CPU/memory optimisations); RSS
unchanged (we are not allocating new big tensors). The 16 GB cgroups
cap applies to any post-optimisation training runs.

\section{Rollback path}

Each pass is a single-module diff. Revert via git revert of the
specific commit. P5 and P6 are independent; P7 depends on P6's
edge-lookup contract; P8 is a separate architectural phase.

\section{Risk anticipation}

\begin{itemize}
\item \emph{Mid-ablation source modifications.} A source modification
  during a running ablation would taint the experimental comparison
  because Python re-imports between configs. \textbf{Strict rule:}
  no source changes until the ablation completes. This plan is
  applied AFTER the 5-config $\times$ 1-seed ablation
  (\texttt{ricci-stim-ablation-2026-05-15.service}) returns its full table.

\item \emph{Diminishing returns.} 296$\times$ already delivered; the
  next 10$\times$ requires P8 (architectural). P5--P7 alone get us
  perhaps 2$\times$ more (14 ms / image / 70 FPS).

\item \emph{Smoke regressions.} The smoke test (30 imgs $\times$
  1 epoch $\times$ Config E) is the canonical sanity gate after
  each pass. If it slows down, revert immediately.

\item \emph{Empty-plan-dir hygiene.} If P5 reveals there is no
  single big-win op in the unaccounted 10 ms and we shelve the
  plan, the plan dir is deleted before commit.
\end{itemize}

\section{Decision tree}

\begin{itemize}
\item \textbf{Config E in the ablation $\geq$ 0.5 mAP50\_proxy:}
  proceed with P5--P7 to get 28 $\to$ 14 ms (70 FPS). P8 deferred
  unless training wall-time becomes the gate.
\item \textbf{Config E in $[0.235, 0.5)$:} architecture is partly
  validated; P5--P7 + a second-seed run becomes the priority over
  raw FPS gains.
\item \textbf{Config E $< 0.235$:} stimulus-driven hypothesis
  falsified; optimization passes are deferred while we
  reconsider the architecture.
\end{itemize}

\section{Bottom line}

Architecturally complete, empirically validated (or not — pending
ablation), and now ready for a focused four-pass per-image
speedup sweep. The plan is paper-only until the in-flight
ablation lands; no source modifications mid-run.

\end{document}
